\section{Déploiement réparti}\label{sec:reparti}
\subsection{Architecture et placements}
Le déploiement utilise le sharding natif de MongoDB : deux shards (\code{:27021},
\code{:27022}) avec chacun son volume, un serveur de configuration (\code{:27020}) et un
routeur \code{mongos} (\code{:27019}). L'API ne change pas, seule \code{MONGODB_URI}
change. Chaque trajet est stocké sur un seul shard, sans copie. Chaque shard est un
replica set à un seul membre, parce que MongoDB l'impose. \code{zones} et \code{imports}
restent sur shard1. Les trois déploiements ont les mêmes index. Le balancer est arrêté
et les plages sont fixées avant l'import : aucune migration ne perturbe les mesures.

\textbf{Placement par zone.} La clé est (zone de départ, date de départ). Un découpage
par borough ne tiendrait pas : Manhattan compte 89\,\% des départs de janvier
(2\,646\,948 sur 2\,964\,624) et Brooklyn 25\,258. Nous classons donc les zones par
nombre de départs sur janvier--mars, et chacune va sur le shard le moins chargé. Le
résultat est de 134 zones sur shard1, 131 sur shard2, et 4\,777\,389 départs de chaque
côté. Les zones très actives de la charge A se retrouvent séparées : 161 et 236 sur
shard1, 132 et 237 sur shard2. L'équilibre tient à six mois (10\,167\,898 et 10\,164\,195
documents) alors qu'il est calculé sur trois. Le mois de F6 et les suivants suivent les mêmes plages de zones.

\textbf{Placement par période (alternative).} La clé est (date de départ, zone) et la
limite est le 1\textsuperscript{er} avril 2024 : janvier--mars sur shard1, le mois de F6
et les suivants sur shard2. Ce placement suit la croissance de l'archive et fait lire par un
seul shard les requêtes limitées à un trimestre. Il stocke 9\,554\,795 et 10\,777\,298 documents.

\begin{figure}[!htbp]\centering\footnotesize
\begin{tikzpicture}[
  every node/.style={align=center},
  box/.style={draw=black!60,rounded corners=2pt,inner sep=4pt,minimum height=9mm},
  link/.style={-{Latex[length=1.6mm]},draw=black!60},
  note/.style={font=\scriptsize,text=black!65}]
\node[box] (client) at (0,0) {Analyste\\(Postman)};
\node[box] (api) at (3.1,0) {API Spring Boot\\F1 à F6};
\node[box] (ingest) at (3.1,1.55) {Ingestion Python\\\code{urbanflow.ingest}};
\node[box] (mongos) at (6.9,0) {\code{mongos}\\\code{:27019}};
\node[box] (config) at (6.9,-1.55) {Configuration\\\code{:27020}};
\node[box,minimum width=6.0cm] (s1) at (12.6,0.95) {\textbf{shard1} \code{:27021} (avec \code{zones}, \code{imports})\\zone : 134 zones, 10\,167\,898 documents\\période : janvier--mars, 9\,554\,795 documents};
\node[box,minimum width=6.0cm] (s2) at (12.6,-0.95) {\textbf{shard2} \code{:27022}\\zone : 131 zones, 10\,164\,195 documents\\période : avril--juin, 10\,777\,298 documents};
\draw[link] (client) -- (api);
\draw[link] (api) -- node[above,note] {requêtes} (mongos);
\draw[link] (api) -- node[right,note] {F6} (ingest);
\draw[link] (ingest.east) -| node[pos=0.25,above,note] {lots de 5\,000 lignes} (mongos.north);
\draw[link] (mongos) -- (config);
\draw[link] (mongos.east) -- ++(0.35,0) |- (s1.west);
\draw[link] (mongos.east) -- ++(0.35,0) |- (s2.west);
\end{tikzpicture}
\caption{Déploiement réparti, avec les documents stockés à six mois. Le déploiement local garde la même API et le même import, avec un seul serveur sur \code{:27018}. Un seul déploiement tourne à la fois, chacun avec ses propres volumes.}\label{fig:architecture}
\end{figure}


\subsection{Ce que fait chaque instance}
\code{explain-requests.py} récupère le pipeline réellement envoyé par l'API et le rejoue
avec \code{explain} sur chaque déploiement \cite{experiences}.

\begin{center}\footnotesize
\begin{tabularx}{\linewidth}{@{}llllX@{}}
\toprule
Requête & Déploiement & Contactés & Documents lus & Sur les shards $\to$ fusion\\
\midrule
Une zone chargée & Local & serveur & 157\,704 & tout sur le serveur\\
(F1, JFK, mars) & Zone & shard2 & 157\,704 & tout sur shard2, sans fusion\\
 & Période & shard1, shard2 & 157\,704 / 0 & \code{$match}, \code{$project} $\to$ \code{$facet}\\
\midrule
Plusieurs zones & Local & serveur & 189\,938 & tout sur le serveur\\
(F2, 5 zones, semestre) & Zone & shard1, shard2 & 88\,485 / 101\,453 & \code{$match}, \code{$project}, \code{$group} $\to$ \code{$group}\\
 & Période & shard1, shard2 & 86\,743 / 103\,195 & idem\\
\midrule
Origine--destination & Local & serveur & 4\,119\,933 & tout sur le serveur\\
(F5, ville, autour & Zone & shard1, shard2 & 2\,062\,737 / 2\,057\,196 & \code{$match}, \code{$project}, \code{$group} $\to$ \code{$group}\\
du 1\textsuperscript{er} avril) & Période & shard1, shard2 & 3\,245\,508 / 874\,425 & idem\\
\bottomrule
\end{tabularx}
\end{center}

Pour F2 et F5, chaque shard regroupe ses documents et n'envoie que des comptes partiels ;
le plan de fusion est \code{$mergeCursors} puis \code{$group}, sur un shard choisi par
MongoDB (\code{anyShard}), et Java combine ensuite les comptes reçus. Sur F2 cinq zones
en placement par zone, les shards émettent 11 Mo par appel et \code{mongos} en rend 4,7 à l'API. \code{$facet}
n'est pas découpé : en placement par période, F1 sur JFK fait sortir 34 Mo par appel de
shard1, contre moins de 0,1 Mo depuis l'unique shard du placement par zone, qui exécute
tout le pipeline sur place.

\subsection{Comparaison automatique}
\code{compare-deployments.py} compare 13 réponses : les huit requêtes de E2 et cinq
requêtes de F1 à F5 qui traversent le 1\textsuperscript{er} avril, plusieurs zones ou les
deux. Les clés JSON sont triées. Les listes gardent l'ordre des services, qui
départagent les ex æquo par identifiants, donc une valeur ou un rang différent fait
échouer le script. Les 13 réponses sont identiques sur les trois déploiements. Une empreinte par fichier (documents, somme des distances, somme des
montants) confirme aussi que les trois bases contiennent les mêmes totaux
\cite{experiences}.

\subsection{Perte d'une instance}
\code{lose-instance.py} arrête un shard, appelle F1 sur la zone 181 à l'arrivée sur le
semestre, dont les trajets sont sur les deux shards, puis F5 sur la ville. Sans shard1 ou
sans shard2, sur les deux placements, l'utilisateur reçoit en 15 à 19 s une erreur 503 :
« MongoDB indisponible ou requête interrompue. Aucun résultat complet ne peut être
fourni. » Le script compare toute réponse 200 à une référence prise sans panne et échoue
si elle diffère.

Shard2 arrêté, une requête limitée à shard1 répond encore 200, à l'identique, parce
qu'elle ne dépend pas de shard2 : zone 1 au départ (placement par zone), zone 181 en
janvier (placement par période). Shard1 arrêté, ces requêtes échouent aussi, parce que
\code{zones} et \code{imports} sont sur ce shard. Après le redémarrage, toutes les
réponses redeviennent identiques à la référence.

Notre déploiement ne récupère donc rien pendant la panne : il n'existe aucune copie, et
les données reviennent seulement avec le shard et son volume. Un système de production
avec des replica sets de plusieurs membres pourrait élire un nouveau primaire et
continuer à répondre ; nous ne l'avons ni déployé ni testé. Nos shards sont aussi deux
conteneurs sur un même PC, dont les pannes ne sont pas indépendantes comme sur des
machines séparées.
